%-----------------------------------------------------------------------
\subsection{Underdetermined Linear Systems (Mathematics)}
\label{sec:instance-linear-system}
%-----------------------------------------------------------------------

The most elementary instance of the Explanation Impossibility arises in
numerical linear algebra \citep{strang2016introduction}.  Consider the system $Ax = b$ with $A = [1\;\;1]$
and $b = 2$.  The solution set $\{x \in \mathbb{R}^2 : x_1 + x_2 = 2\}$ is a
line, not a point: both $x^{(1)} = (1,1)$ and $x^{(2)} = (0,2)$ satisfy
$Ax = b$.  Every STEM student encounters this situation in their first
linear algebra course.

\paragraph{The explanation system.}
Define $\mathcal{S}_{\mathrm{lin}} = (\Params, \sH, Y, \mathsf{obs},
\mathsf{exp}, \incomp)$ as follows:
\begin{itemize}
  \item $\Params = \sH = \mathbb{Z}^2$ (solution vectors);
  \item $Y = \mathbb{Z}$ (the observable: $Ax$);
  \item $\mathsf{obs}(v) = v_1 + v_2$ (dot product with $A = [1\;\;1]$);
  \item $\mathsf{exp} = \mathrm{id}$ (the solution itself is the ``explanation'');
  \item $\incomp(v, v') \iff v \ne v'$.
\end{itemize}

\paragraph{The Rashomon property.}
The witnesses $v^{(1)} = (1,1)$ and $v^{(2)} = (0,2)$ satisfy
$\mathsf{obs}(v^{(1)}) = 1 + 1 = 2 = 0 + 2 = \mathsf{obs}(v^{(2)})$
and $v^{(1)} \ne v^{(2)}$.  The Rashomon property is constructively witnessed
with zero axioms.

\paragraph{The impossibility.}
By Theorem~\ref{thm:universal}, no method of choosing a solution to an
underdetermined linear system can simultaneously be:
\begin{enumerate}
  \item \textbf{Faithful}: the chosen solution agrees with the actual solution;
  \item \textbf{Stable}: all solutions with the same image receive the same
    choice; and
  \item \textbf{Decisive}: the choice distinguishes every pair of distinct
    solutions.
\end{enumerate}

\paragraph{Empirical illustration.}
Four solvers (pseudoinverse, LSQR, Tikhonov, random null-space) applied to 100
underdetermined systems show mean pairwise RMSD of 0.091, compared to 0.049 for
fully determined controls (Mann-Whitney $p = 6.3 \times 10^{-9}$).
Disagreement is elevated across all null space dimensions $d > 0$.
\emph{Note:} RMSD does not increase monotonically with $d$ because three of
the four non-random solvers converge to near-minimum-norm solutions; the
disagreement is therefore dominated by the single random null-space projection
solver.  The key finding is the positive/zero separation ($d > 0$ vs.\ $d = 0$),
not a smooth scaling relationship.

\paragraph{Resolution: minimum-norm solutions.}
The classical resolution is the Moore--Penrose pseudoinverse
$x^+ = A^+(Ax)$, which selects the minimum-norm element of the solution
set.  This is precisely a $G$-invariant projection: the symmetry group
$G$ acts by translations along the null space of $A$, and $x^+$ is the unique
element of each fiber that is invariant under the orthogonal complement of
$\ker(A)$.  The pseudoinverse sacrifices decisiveness (it cannot distinguish
solutions within the same fiber) to achieve stability and faithfulness
in expectation.

\paragraph{Lean formalization.}
The Lean~4 types are \texttt{Vec2} (a structure with fields \texttt{x y : Int}),
\texttt{dotA} (the observation map $v \mapsto v.x + v.y$), and
\texttt{linearSystem : ExplanationSystem Vec2 Vec2 Int} with
\texttt{observe := dotA}, \texttt{explain := id}, and
\texttt{incompatible := ($\ne$)}.
The impossibility is \texttt{linear\_system\_impossibility} in
\texttt{LinearSystem.lean}, derived with zero axioms.
